\documentclass[a4paper,answers]{exam} % Pour la version avec les réponses, utiliser: \documentclass[answers]{exam}
\usepackage[utf8]{inputenc}
\usepackage[T1]{fontenc}
\usepackage{amsmath, amssymb}
\usepackage{siunitx}

\newcommand{\pts}[1]{\textbf{#1 pts}}

\extrawidth{1in}

%\firstpageheader{
%\large{Évaluation - Projet 2 : Modélisation moléculaire}

\pagestyle{headandfoot}
\runningheadrule
\firstpageheader{PHY 4268}{Évaluation Projet 2 - Modélisation moléculaire}{21 Mars 2025}
\runningheader{PHY 4268}
{Évaluation Projet 2 - Modélisation moléculaire, Page \thepage\ of \numpages}
{21 Mars 2025}
\firstpagefooter{}{}{}
\runningfooter{}{}{}

\begin{document}

\begin{itemize}
    \item Durée : 20 minutes
\end{itemize}

\vspace{0.2cm}

\begin{questions}

\question[6] \textbf{Hamiltonien moléculaire électronique.} Nommez et décrivez brièvement (en une phrase) les \textbf{trois termes de potentiel} qui composent l'Hamiltonien moléculaire électronique ($\mathtt{H}_{\rm elec}$).

\begin{solution}
\begin{itemize}
    \item $\mathtt{U}_{en}= - \sum_{i,I}\frac{Z_I}{|\mathbf{r}_i-\mathbf{R}_I|}$ (Énergie potentielle d'attraction électron-noyau) \hfill \pts{1}
        \begin{itemize}
            \item Interaction attractive Coulombienne entre électrons et noyaux. \hfill \pts{1}
        \end{itemize}
    \item $\mathtt{U}_{ee}=+\sum_{i,j>i}\frac{1}{r_{ij}}$ (Énergie potentielle de répulsion électron-électron) \hfill \pts{1}
        \begin{itemize}
            \item Interaction répulsive Coulombienne entre électrons. \hfill \pts{1}
        \end{itemize}
    \item $\mathtt{U}_{nn}=+\sum_{I,J>I}\frac{Z_IZ_J}{R_{IJ}}$ (Énergie potentielle de répulsion noyau-noyau) \hfill \pts{1}
        \begin{itemize}
            \item Interaction répulsive Coulombienne entre noyaux.  \hfill \pts{1}
        \end{itemize}
\end{itemize}
\end{solution}

\question[5] \textbf{Approximation de Born-Oppenheimer.} Expliquez en \textbf{termes simples} (2-3 phrases maximum) le principe de l'Approximation de Born-Oppenheimer. Quel est l'\textbf{avantage principal} de cette approximation pour la modélisation moléculaire ?

\begin{solution}
\begin{itemize}
    \item \textbf{Principe BO} : Séparation des mouvements noyaux/électrons due à la différence de masse. Noyaux considérés fixes pour le mouvement électronique. \hfill \pts{3}
    \item \textbf{Avantage principal} : Simplification drastique des calculs, rendant la modélisation possible.\hfill\pts{2}
\end{itemize}
\end{solution}

\question[5] \textbf{Méthode de Hartree-Fock et champ moyen ou mean-field.}
\begin{parts}
    \part[2] Qu'est-ce que l'approximation du "\textbf{champ moyen}" dans la méthode de Hartree-Fock ? (1 phrase)
    \begin{solution}
Champ moyen : Chaque électron se déplace dans un potentiel moyen créé par l'ensemble des autres électrons (interaction électron-électron moyennée). \hfill \pts{2}
    \end{solution}

    \part[3] Pourquoi la méthode de Hartree-Fock est-elle qualifiée de procédure "\textbf{auto-cohérente (SCF)}"? (1 phrase)
    \begin{solution}
Auto-cohérence (SCF) : La procédure est itérative car les orbitales (et donc le champ moyen) sont recalculées et mises à jour jusqu'à convergence. \hfill \pts{3}
    \end{solution}
\end{parts}

\question[5] \textbf{Choix de la méthode Hartree-Fock.} Dans quel type de situation choisiriez-vous la méthode Hartree-Fock \textbf{UHF} plutôt que la méthode \textbf{UHF} ? \textbf{Justifiez brièvement} votre réponse en considérant les types de systèmes moléculaires et leurs propriétés électroniques.

\begin{solution}
\begin{itemize}
    \item UHF plutôt que RHF : Pour les systèmes à couches ouvertes (radicaux, systèmes magnétiques) ou lorsque RHF n'est pas adapté. \hfill \pts{3}
    \item Justification : UHF permet de traiter les systèmes où les spins alpha et beta ne sont pas contraints d'occuper les mêmes orbitales spatiales, ce qui est nécessaire pour décrire correctement ces systèmes. \hfill \pts{2}
\end{itemize}
\end{solution}


\end{questions}

\end{document} 
